\documentclass[a4paper,12pt]{article}

\usepackage[utf8]{inputenc}
\usepackage[T5]{fontenc}
\usepackage[english,vietnamese]{babel}
\usepackage{mathptmx}
\usepackage{geometry}
\usepackage{enumitem}
\usepackage{xcolor}
\usepackage{hyperref}

\geometry{margin=2.2cm}
\AtBeginDocument{\selectlanguage{vietnamese}}
\setlist[itemize]{leftmargin=1.2cm,itemsep=2pt,topsep=2pt}
\setlist[enumerate]{leftmargin=1.2cm,itemsep=2pt,topsep=2pt}
\hypersetup{colorlinks=true,linkcolor=black,urlcolor=blue}

\newcommand{\slide}[2]{%
  \section*{Slide #1: #2}
  \addcontentsline{toc}{section}{Slide #1: #2}
}
\newcommand{\talk}{\textbf{Lời nói đề xuất.} }
\newcommand{\note}{\textbf{Ý cần nhấn mạnh.} }

\title{\textbf{Speaker Script}\\Ecopark Bicycle Parking -- Use-Case Analysis}
\author{Nhóm dự án hệ thống thuê xe đạp Ecopark}
\date{Tháng 5 năm 2026}

\begin{document}
\maketitle
\tableofcontents
\newpage

\slide{1}{Ecopark Bicycle Parking}
\talk
Chào thầy/cô và các bạn. Nhóm em trình bày phần phân tích use case cho hệ thống
\textit{Ecopark Bicycle Parking}. Đây là hệ thống website hỗ trợ thuê và trả xe
đạp tại các điểm tập kết trong Ecopark. Luồng nghiệp vụ chính đi từ tạo tài
khoản, tìm trạm và chọn xe, đặt cọc nhận xe, sử dụng xe, trả xe và nhận vé
thanh toán cuối cùng.

\note
Nói ngắn gọn phạm vi: không phải app gọi xe chung chung, mà là hệ thống thuê-trả
xe đạp có quản lý trạm, xe, giấy tờ định danh, đặt cọc và lịch sử thuê.

\slide{2}{Project Context and Rules}
\talk
Theo đề bài, hệ thống cần quản lý nhiều điểm tập kết xe trong Ecopark. Mỗi xe có
mã định danh duy nhất toàn hệ thống và có trạng thái như sẵn sàng, đang được
thuê hoặc hỏng cần sửa. Người dùng lần đầu thuê xe phải tạo tài khoản với email,
số điện thoại, địa chỉ và CCCD hoặc giấy tờ định danh. Khi nhận xe, người dùng
gặp nhân viên tại trạm để đặt cọc 200 nghìn và để lại giấy tờ tùy thân. Khi trả
xe, hệ thống tính tiền theo gói giờ, cộng phí trễ nếu có, và lưu vé vào lịch sử.

\note
Nhấn mạnh quy tắc giảm giá 40\% cho cư dân có thẻ cư dân hợp lệ. Đồng thời nói
rõ front-end, back-end, database và bảo vệ cổng không được xem là actor trong
use-case diagram.

\slide{3}{Actor Model}
\talk
Ở mô hình actor, nhóm em xác định năm actor chính. Customer/User là người thuê
xe thông thường. Resident Customer là actor chuyên biệt kế thừa Customer/User,
vì cư dân vẫn có toàn bộ chức năng của người thuê thông thường nhưng có thêm
nhánh xác minh thẻ cư dân để được giảm giá. Station Staff hoặc Manager là nhân
viên tại bãi, phụ trách kiểm tra giấy tờ, nhận cọc, giao xe và nhận xe trả.
Admin hoặc Operator quản lý dữ liệu toàn hệ thống như bãi xe, xe, chính sách giá
và báo cáo. GPS/Map Service là dịch vụ ngoài hỗ trợ tìm trạm gần nhất.

\note
Trả lời trước câu hỏi hay gặp: Admin không tự làm việc xác minh tài khoản cho
tất cả user, nhưng có thể tham gia manual review khi hồ sơ đáng ngờ hoặc thẻ cư
dân cần rà soát.

\slide{4}{Use-Case Overview}
\talk
Slide này là bức tranh tổng quan của hệ thống. Nhóm em chọn năm use case quan
trọng nhất vì chúng bao phủ toàn bộ vòng đời thuê xe. UC001 xử lý đăng ký và xác
minh tài khoản. UC002 giúp người dùng tìm trạm, xem xe và gửi yêu cầu thuê.
UC003 là nghiệp vụ nhận xe tại trạm, gồm kiểm tra giấy tờ, đặt cọc và bàn giao
xe. UC004 là nghiệp vụ trả xe, tính phí và xuất vé. UC005 dành cho vận hành:
quản lý xe, bãi xe, trạng thái và báo cáo.

\note
Nói rõ đây là use case cấp nghiệp vụ, không phải từng button hoặc từng bước UI.
Các actor chỉ nối với các use case mà họ thật sự tương tác.

\slide{5}{Relational Schema Summary}
\talk
Về dữ liệu, nhóm em dùng các nhóm bảng chính. Nhóm tài khoản gồm USERS,
CUSTOMER\_PROFILES và RESIDENT\_CARDS để biểu diễn cả người dùng thường và cư
dân. Nhóm đội xe gồm STATIONS, BIKE\_TYPES, BIKES và BIKE\_STATUS\_LOGS. Nhóm
yêu cầu thuê gồm RENTAL\_REQUESTS. Khi nhận xe thành công, hệ thống tạo RENTALS
và RENTAL\_DEPOSITS. Khi trả xe, hệ thống tạo RENTAL\_TICKETS. Cuối cùng,
PRICING\_POLICIES lưu các mức giá theo giờ, phí trễ và tỷ lệ giảm giá cư dân.

\note
Nhấn mạnh schema không dùng ví điện tử hay gói ngày/tháng vì đề mới không yêu
cầu; nhóm đã điều chỉnh theo đề bài hiện tại.

\slide{6}{Include / Extend Notes for Diagram}
\talk
Trong sơ đồ chi tiết, nhóm em dùng include và extend ở mức vừa đủ. Include được
dùng khi một use case con luôn là một phần bắt buộc của use case chính, ví dụ
xác minh giấy tờ và nhận cọc trong nghiệp vụ bàn giao xe. Extend được dùng khi
một nhánh chỉ xảy ra trong điều kiện nhất định, ví dụ nộp thẻ cư dân khi người
dùng là cư dân, hoặc áp dụng phí trễ khi xe được trả muộn.

\note
Khi bị hỏi về ký hiệu, trả lời: mũi tên include/extend chỉ nối giữa use case với
use case, không nối actor với actor hoặc actor với database.

\slide{7}{UC001 Main Flow: Register and Verify Customer Account}
\talk
UC001 bắt đầu khi người dùng đăng ký hoặc đăng nhập. Hệ thống yêu cầu các thông
tin bắt buộc gồm email, số điện thoại, địa chỉ, CCCD hoặc giấy tờ định danh. Sau
khi người dùng gửi thông tin, hệ thống kiểm tra định dạng, trùng email hoặc số
điện thoại, và độ mạnh mật khẩu. Nếu hợp lệ, hệ thống tạo tài khoản và hồ sơ
khách hàng. Trước lần thuê đầu tiên, người dùng phải có thông tin định danh hợp
lệ để được đánh dấu đủ điều kiện thuê. Nếu là cư dân, họ có thể nộp địa chỉ
Ecopark và thẻ cư dân để được xét giảm giá.

\note
Nhấn mạnh mượn xe là tài sản có giá trị, nên phần đăng ký/đăng nhập phải chặt
hơn một ứng dụng xem thông tin thông thường.

\slide{8}{UC001 Alternative Flows}
\talk
Các luồng thay thế của UC001 chủ yếu xử lý dữ liệu không hợp lệ hoặc chưa đủ.
Nếu thiếu email, số điện thoại, địa chỉ hoặc giấy tờ định danh, hệ thống không
cho thuê và yêu cầu bổ sung. Nếu email hoặc số điện thoại đã tồn tại, hệ thống
chặn đăng ký mới và gợi ý đăng nhập hoặc khôi phục tài khoản. Nếu giấy tờ bị sai
hoặc bị khóa, người dùng không đủ điều kiện thuê. Với cư dân, nếu không có thẻ
cư dân thì vẫn được thuê như khách thường, nhưng không được giảm 40\%. Nếu thẻ
cư dân đáng ngờ, hồ sơ chuyển sang trạng thái chờ Admin/Operator rà soát.

\note
Đây là nơi giải thích vì sao Resident Customer kế thừa Customer/User: cư dân chỉ
thêm nhánh giảm giá, không phải một hệ thống thuê xe riêng.

\slide{9}{UC002 Main Flow: Find Station and Select Bicycle}
\talk
UC002 bắt đầu khi người dùng mở màn hình tìm trạm hoặc thuê xe. Hệ thống lấy vị
trí hiện tại hoặc khu vực tìm kiếm thủ công. GPS/Map Service trả về vị trí hoặc
khu vực, sau đó hệ thống liệt kê các trạm gần nhất cùng số lượng xe còn trống
theo loại. Người dùng mở chi tiết trạm để xem mã xe, loại xe và trạng thái từng
xe. Người dùng có thể lọc theo loại xe, chọn xe còn sẵn sàng, xác nhận thời
lượng dự kiến và gửi yêu cầu thuê. Kết quả là một RentalRequest chờ xử lý tại
trạm.

\note
Nói rõ “tìm trạm gần nhất”, “xem tồn kho xe”, “lọc loại xe” là các use case con
hoặc nhánh trong mục tiêu lớn hơn, không nên tách thành các actor.

\slide{10}{UC002 Alternative Flows}
\talk
Nếu người dùng từ chối cấp vị trí, hệ thống cho tìm trạm thủ công. Nếu GPS/Map
Service không hoạt động, hệ thống cũng quay về tìm kiếm thủ công hoặc dùng khu
vực gần nhất đã biết. Nếu không có trạm gần, hệ thống yêu cầu mở rộng phạm vi.
Nếu loại xe mong muốn hết, người dùng có thể chọn loại khác hoặc trạm khác. Nếu
xe vừa được người khác thuê hoặc chuyển sang trạng thái hỏng, hệ thống làm mới
tồn kho và yêu cầu chọn lại. Nếu tài khoản thiếu thông tin định danh, hệ thống
chuyển người dùng về UC001.

\note
Nhấn mạnh alternative flow không chỉ là lỗi kỹ thuật, mà còn là tình huống nghiệp
vụ: hết xe, xe hỏng, trạm không hoạt động, tài khoản chưa đủ điều kiện.

\slide{11}{UC003 Main Flow: Process Bicycle Rental and Deposit}
\talk
UC003 là bước nhận xe tại trạm. Người dùng đưa mã yêu cầu thuê cho Station
Staff/Manager. Nhân viên tìm yêu cầu đang chờ xử lý trong hệ thống. Hệ thống
hiển thị thông tin khách hàng, xe, trạm và thời lượng dự kiến. Nhân viên kiểm
tra CCCD hoặc giấy tờ tùy thân và xác nhận giấy tờ khớp với tài khoản. Sau đó
nhân viên nhận cọc 200 nghìn, ghi nhận loại giấy tờ đang giữ, kiểm tra lại trạng
thái trạm và xe. Nếu mọi thứ hợp lệ, nhân viên giao đúng xe cho khách, hệ thống
tạo Rental, ghi nhận tiền cọc, thời gian bắt đầu và chuyển xe sang trạng thái
đang thuê.

\note
Customer chỉ báo cáo/yêu cầu thuê; nhân viên mới xử lý bàn giao và đặt cọc. Đây
là điểm phân quyền quan trọng.

\slide{12}{UC003 Alternative Flows}
\talk
Nếu mã yêu cầu thuê không tồn tại hoặc đã hết hạn, nhân viên từ chối bàn giao và
yêu cầu người dùng tạo yêu cầu mới. Nếu yêu cầu thuộc trạm khác, trạm hiện tại
không được xử lý. Nếu người dùng quên giấy tờ, giấy tờ không khớp, không đủ cọc
hoặc từ chối để lại giấy tờ, xe không được bàn giao. Nếu xe đã hỏng, không còn ở
trạm hoặc dữ liệu vị trí bị lệch, hệ thống chặn bàn giao và ghi nhận để vận hành
xử lý. Nếu lưu rental hoặc deposit thất bại, xe không được đánh dấu là đang thuê.

\note
Kết luận ngắn: UC003 chỉ hoàn tất khi cả bốn điều kiện đều qua: request hợp lệ,
identity hợp lệ, cọc hợp lệ và xe sẵn sàng.

\slide{13}{UC004 Main Flow: Return Bicycle and Issue Ticket}
\talk
UC004 xảy ra khi người dùng mang xe tới bất kỳ trạm nào trong Ecopark để trả.
Nhân viên tìm rental theo mã thuê, khách hàng hoặc mã xe. Hệ thống hiển thị
thông tin chuyến thuê, trạm nhận xe ban đầu, thời gian bắt đầu, thời lượng dự
kiến và thông tin cọc. Nhân viên kiểm tra tình trạng xe, xác nhận trạm trả và
thời gian trả. Hệ thống tính thời lượng thực tế, tính phí cơ bản theo các mức
50 nghìn một giờ, 70 nghìn hai giờ hoặc 100 nghìn ba giờ. Nếu khách là cư dân
đủ điều kiện, hệ thống giảm 40\%. Nếu trả muộn, hệ thống cộng 30 nghìn cho mỗi
30 phút. Cuối cùng hệ thống xuất vé, lưu lịch sử và cập nhật vị trí/trạng thái
xe.

\note
Nói rõ đề bài cho phép trả xe ở bất kỳ trạm nào trong Ecopark, không bắt buộc
trả đúng trạm đã nhận.

\slide{14}{UC004 Alternative Flows}
\talk
Nếu trạm trả khác trạm nhận, hệ thống vẫn cho trả và cập nhật vị trí xe mới. Nếu
không tìm thấy rental đang chạy, hệ thống chuyển sang rà soát thủ công. Nếu thông
tin cọc bị lệch, cần đánh dấu để đối soát trước khi đóng. Nếu xe trả về bị hỏng,
hệ thống chuyển trạng thái xe sang broken hoặc needs repair và không cho thuê
tiếp. Nếu nhân viên nhập sai thời gian hoặc trạm trả, hệ thống cho sửa trước khi
phát hành vé. Nếu khách tranh chấp phí trễ hoặc thời gian, vé vẫn được lưu nhưng
đánh dấu cần Admin/Operator rà soát.

\note
Nhấn mạnh mọi trường hợp trả xe đều cần kết quả có thể audit: vé, lịch sử, trạng
thái xe và lý do nếu cần review.

\slide{15}{UC005 Main Flow: Manage Bikes, Stations and Rental Reports}
\talk
UC005 dành cho vận hành hệ thống. Admin/Operator hoặc Staff mở màn hình quản lý
vận hành và chọn quản lý trạm, quản lý xe, cập nhật trạng thái, định vị xe hoặc
xem báo cáo. Với trạm, actor có thể tạo hoặc cập nhật tên, địa chỉ, tọa độ, sức
chứa và trạng thái hoạt động. Với xe, actor quản lý mã xe duy nhất, loại xe, vị
trí hiện tại và trạng thái. Mọi thay đổi trạng thái hoặc vị trí đều được hệ
thống ghi log với người thực hiện và lý do. Với báo cáo, Admin chọn khoảng thời
gian, trạm, loại xe hoặc mã xe, sau đó hệ thống tổng hợp lượt thuê, doanh thu,
phí trễ, lượt trả và trạng thái xe.

\note
Nói rõ Staff có thể tham gia vận hành ở phạm vi trạm, còn Admin/Operator có
quyền toàn hệ thống và xem báo cáo tổng hợp.

\slide{16}{UC005 Alternative Flows}
\talk
Nếu nhân viên trạm không có quyền toàn hệ thống, hệ thống giới hạn dữ liệu và
hành động trong trạm được phân công. Nếu tạo trạm thiếu địa chỉ, tọa độ hoặc sức
chứa, hệ thống từ chối lưu. Nếu sức chứa mới nhỏ hơn số xe đang có tại trạm, hệ
thống không cho cập nhật. Nếu mã xe bị trùng, hệ thống từ chối tạo hoặc sửa xe.
Nếu xe đang được thuê, hệ thống không cho đổi vị trí thủ công cho tới khi rental
kết thúc. Nếu xe hỏng, xe bị loại khỏi danh sách có thể thuê; muốn mở lại phải
có ghi chú kiểm tra hoặc sửa chữa. Nếu báo cáo không có dữ liệu, hệ thống hiển
thị rỗng thay vì báo lỗi.

\note
Nhấn mạnh UC005 là use case vận hành trung tâm, giúp đảm bảo dữ liệu xe và báo
cáo không bị sai lệch.

\slide{17}{Presentation Checklist}
\talk
Để kết thúc, nhóm em sẽ trình bày theo thứ tự: bắt đầu từ bối cảnh và actor,
sau đó đi vào sơ đồ use case tổng quan, giải thích schema quan hệ và lần lượt
trình bày năm use case chính. Với mỗi use case, nhóm nêu actor, mục tiêu, kết
quả quan sát được, main flow và alternative flow. Nhóm cũng đã tránh các lỗi mô
hình hóa phổ biến như xem front-end, back-end hoặc database là actor, hoặc đặt
tên use case như một bước UI quá nhỏ.

\note
Khi bị hỏi chi tiết, mở file specification A4 để chỉ bảng đầy đủ của từng UC:
preconditions, postconditions, input, output, main flow, alternative flow và
acceptance rules.

\end{document}
